%============================================================
% RAPPORT DE PROJET DE FIN D'ÉTUDES (PFE)
% Filière : Ingénierie en Informatique
% Projet  : VulnOps – Plateforme DevSecOps Intelligente
%============================================================

\documentclass[12pt, a4paper, oneside]{report}

%------------------------------------------------------------
% PACKAGES
%------------------------------------------------------------
\usepackage[utf8]{inputenc}
\usepackage[T1]{fontenc}
\usepackage[french]{babel}
\usepackage{geometry}
\geometry{top=2.5cm, bottom=2.5cm, left=3cm, right=2.5cm}

\usepackage{graphicx}
\usepackage{xcolor}
\usepackage{hyperref}
\usepackage{amsmath}
\usepackage{amssymb}
\usepackage{listings}
\usepackage{booktabs}
\usepackage{longtable}
\usepackage{array}
\usepackage{tabularx}
\usepackage{multirow}
\usepackage{float}
\usepackage{caption}
\usepackage{subcaption}
\usepackage{fancyhdr}
\usepackage{titlesec}
\usepackage{tocloft}
\usepackage{setspace}
\usepackage{parskip}
\usepackage{enumitem}
\usepackage{lmodern}
\usepackage{microtype}
\usepackage{pgfgantt}
\usepackage{rotating}
\usepackage{pdfpages}
\usepackage{cite}
\usepackage{url}
\usepackage{tikz}
\usetikzlibrary{shapes.geometric, arrows.meta, positioning, fit, backgrounds}

%------------------------------------------------------------
% COULEURS
%------------------------------------------------------------
\definecolor{primaryblue}{RGB}{0, 84, 166}
\definecolor{secondaryblue}{RGB}{0, 150, 214}
\definecolor{darkgray}{RGB}{50, 50, 50}
\definecolor{lightgray}{RGB}{245, 245, 245}
\definecolor{codegreen}{RGB}{0, 128, 0}
\definecolor{codepurple}{RGB}{128, 0, 128}
\definecolor{codeblue}{RGB}{0, 0, 200}
\definecolor{codered}{RGB}{160, 0, 0}

%------------------------------------------------------------
% CONFIGURATION HYPERREF
%------------------------------------------------------------
\hypersetup{
    colorlinks=true,
    linkcolor=primaryblue,
    citecolor=primaryblue,
    urlcolor=secondaryblue,
    pdftitle={Rapport PFE -- VulnOps},
    pdfauthor={Firas Halouani},
    pdfsubject={Plateforme DevSecOps Intelligente},
}

%------------------------------------------------------------
% CONFIGURATION LISTINGS (CODE)
%------------------------------------------------------------
\lstset{
    basicstyle=\ttfamily\footnotesize,
    keywordstyle=\color{codeblue}\bfseries,
    commentstyle=\color{codegreen}\itshape,
    stringstyle=\color{codered},
    numberstyle=\tiny\color{darkgray},
    numbers=left,
    numbersep=5pt,
    backgroundcolor=\color{lightgray},
    breaklines=true,
    breakatwhitespace=true,
    frame=single,
    framerule=0.4pt,
    tabsize=4,
    captionpos=b,
    showstringspaces=false,
    extendedchars=true,
}

%------------------------------------------------------------
% EN-TÊTES ET PIEDS DE PAGE
%------------------------------------------------------------
\pagestyle{fancy}
\fancyhf{}
\fancyhead[L]{\small\leftmark}
\fancyhead[R]{\small\textcolor{primaryblue}{VulnOps -- Rapport PFE}}
\fancyfoot[C]{\thepage}
\renewcommand{\headrulewidth}{0.4pt}
\renewcommand{\footrulewidth}{0pt}

%------------------------------------------------------------
% STYLE DES TITRES DE CHAPITRES
%------------------------------------------------------------
\titleformat{\chapter}[display]
  {\normalfont\huge\bfseries\color{primaryblue}}
  {\chaptertitlename\ \thechapter}
  {20pt}
  {\Huge}

\titleformat{\section}
  {\normalfont\Large\bfseries\color{primaryblue}}
  {\thesection}
  {1em}{}

\titleformat{\subsection}
  {\normalfont\large\bfseries\color{darkgray}}
  {\thesubsection}
  {1em}{}

\titleformat{\subsubsection}
  {\normalfont\normalsize\bfseries}
  {\thesubsubsection}
  {1em}{}

%------------------------------------------------------------
% INTERLIGNE
%------------------------------------------------------------
\setstretch{1.25}

%============================================================
% DÉBUT DU DOCUMENT
%============================================================
\begin{document}

%------------------------------------------------------------
% PAGE DE GARDE
%------------------------------------------------------------
\begin{titlepage}
    \centering
    \vspace*{1cm}

    {\large\textbf{République Tunisienne}}\\[0.3cm]
    {\large Ministère de l'Enseignement Supérieur et de la Recherche Scientifique}\\[0.5cm]
    \hrule height 1.5pt
    \vspace{0.3cm}
    {\Large\textbf{Université [Votre Université]}}\\[0.2cm]
    {\large Institut/École Supérieure [Votre Établissement]}\\[0.3cm]
    \hrule height 1.5pt

    \vspace{1.5cm}

    % Logo établissement
    % \includegraphics[width=4cm]{figures/logo_etablissement.png}\\[1cm]

    \vspace{0.5cm}
    {\LARGE\textbf{\textcolor{primaryblue}{Rapport de Projet de Fin d'Études}}}\\[0.3cm]
    {\large Pour l'obtention du diplôme d'Ingénieur en Informatique}\\[0.3cm]
    {\large Spécialité : Génie Logiciel / Sécurité Informatique}

    \vspace{1cm}
    \hrule
    \vspace{0.5cm}
    {\Huge\bfseries\textcolor{primaryblue}{VulnOps}}\\[0.4cm]
    {\Large\textbf{Plateforme DevSecOps Intelligente pour}}\\[0.2cm]
    {\Large\textbf{l'Analyse Automatisée des Vulnérabilités}}
    \vspace{0.5cm}
    \hrule

    \vspace{1.5cm}
    \begin{tabular}{ll}
        \textbf{Réalisé par :} & Firas Halouani \\[0.3cm]
        \textbf{Encadrant académique :} & Dr. Selma Belgacem \\[0.3cm]
        \textbf{Encadrant professionnel :} & [Nom de l'encadrant] \\[0.3cm]
        \textbf{Organisme d'accueil :} & [Nom de l'entreprise] \\[0.3cm]
        \textbf{Année universitaire :} & 2025 -- 2026 \\
    \end{tabular}

    \vfill
    {\large\textbf{Période de stage :} [Date début] -- [Date fin]}
\end{titlepage}

%------------------------------------------------------------
% DÉDICACE (page non numérotée)
%------------------------------------------------------------
\pagestyle{empty}
\cleardoublepage
\vspace*{5cm}
\begin{center}
    \textit{\Large\textcolor{primaryblue}{Dédicace}}
\end{center}
\vspace{1cm}
\begin{flushright}
    \textit{Je dédie ce modeste travail à mes chers parents,}\\[0.3cm]
    \textit{pour leur soutien indéfectible, leur amour et leurs sacrifices.}\\[0.3cm]
    \textit{À mes frères et sœurs, pour leurs encouragements constants.}\\[0.3cm]
    \textit{À tous mes amis et collègues qui m'ont accompagné}\\[0.3cm]
    \textit{tout au long de ce parcours.}
\end{flushright}
\newpage

%------------------------------------------------------------
% REMERCIEMENTS
%------------------------------------------------------------
\chapter*{Remerciements}
\addcontentsline{toc}{chapter}{Remerciements}
\pagestyle{empty}

Nous tenons, en premier lieu, à exprimer notre profonde gratitude envers
\textbf{Dr. Selma Belgacem}, notre encadrante académique, pour son suivi attentif,
ses conseils avisés et sa disponibilité tout au long de ce projet. Ses orientations
ont été déterminantes pour la réussite de ce travail.

Nous exprimons également nos sincères remerciements à l'ensemble du corps enseignant
de \textbf{[Nom de l'établissement]} pour la qualité de la formation dispensée durant
nos années d'études.

Nos remerciements les plus chaleureux vont aussi à notre encadrant professionnel,
\textbf{[Nom de l'encadrant professionnel]}, ainsi qu'à toute l'équipe de
\textbf{[Nom de l'organisme d'accueil]}, pour leur accueil bienveillant et leur
précieux soutien lors de notre stage.

Nous remercions enfin les membres du jury qui ont bien voulu consacrer de leur
temps à l'évaluation de ce rapport.

\newpage

%------------------------------------------------------------
% TABLE DES MATIÈRES
%------------------------------------------------------------
\pagestyle{plain}
\tableofcontents
\newpage

%------------------------------------------------------------
% LISTE DES FIGURES
%------------------------------------------------------------
\listoffigures
\addcontentsline{toc}{chapter}{Liste des figures}
\newpage

%------------------------------------------------------------
% LISTE DES TABLEAUX
%------------------------------------------------------------
\listoftables
\addcontentsline{toc}{chapter}{Liste des tableaux}
\newpage

%------------------------------------------------------------
% LISTE DES ABRÉVIATIONS
%------------------------------------------------------------
\chapter*{Liste des abréviations}
\addcontentsline{toc}{chapter}{Liste des abréviations}

\begin{tabular}{ll}
    \textbf{API}        & Application Programming Interface \\
    \textbf{CI/CD}      & Continuous Integration / Continuous Delivery \\
    \textbf{CORS}       & Cross-Origin Resource Sharing \\
    \textbf{DAST}       & Dynamic Application Security Testing \\
    \textbf{DevSecOps}  & Development, Security and Operations \\
    \textbf{DRF}        & Django REST Framework \\
    \textbf{IHM}        & Interface Homme-Machine \\
    \textbf{JSON}       & JavaScript Object Notation \\
    \textbf{JWT}        & JSON Web Token \\
    \textbf{LLM}        & Large Language Model \\
    \textbf{MVC}        & Model-View-Controller \\
    \textbf{MVT}        & Model-View-Template \\
    \textbf{OAuth}      & Open Authorization \\
    \textbf{OWASP}      & Open Web Application Security Project \\
    \textbf{PFE}        & Projet de Fin d'Études \\
    \textbf{RAG}        & Retrieval-Augmented Generation \\
    \textbf{REST}       & Representational State Transfer \\
    \textbf{SAST}       & Static Application Security Testing \\
    \textbf{SCA}        & Software Composition Analysis \\
    \textbf{SQL}        & Structured Query Language \\
    \textbf{UI}         & User Interface \\
    \textbf{URL}        & Uniform Resource Locator \\
    \textbf{UML}        & Unified Modeling Language \\
    \textbf{VulnOps}    & Vulnerability Operations \\
\end{tabular}
\newpage

%============================================================
% NUMÉROTATION DES PAGES : démarre à l'introduction générale
%============================================================
\pagestyle{fancy}
\setcounter{page}{1}

%============================================================
% INTRODUCTION GÉNÉRALE
%============================================================
\chapter*{Introduction Générale}
\addcontentsline{toc}{chapter}{Introduction Générale}
\markboth{Introduction Générale}{Introduction Générale}

Dans un contexte où la transformation numérique s'accélère et où les applications
logicielles jouent un rôle central dans tous les secteurs d'activité, la sécurité
informatique est devenue une préoccupation majeure pour les organisations. Le nombre
croissant de cyberattaques et de vulnérabilités exploitées souligne la nécessité
d'intégrer la sécurité dès les premières phases du cycle de développement logiciel.
Les approches traditionnelles, qui traitent la sécurité comme une étape finale du
processus de développement, se révèlent insuffisantes face à la complexité et à la
rapidité d'évolution des menaces actuelles. C'est dans ce cadre que le paradigme
\textbf{DevSecOps} (Development, Security et Opérations) émerge comme une réponse
adaptée, en intégrant les pratiques de sécurité directement au cœur du pipeline
de développement et de déploiement \cite{devsecopsgartner}.

Le présent rapport décrit le travail réalisé dans le cadre de notre stage de fin d’études au sein de \textbf{Mobelite}. Ce stage a abouti à la conception et au développement de VulnOps, une plateforme DevSecOps intelligente dédiée à l’analyse automatisée des vulnérabilités dans les projets logiciels. VulnOps exploite une combinaison d’outils d’analyse statique (SAST), de composition logicielle (SCA), d’analyse des conteneurs (Container Scan) et d’analyse dynamique (DAST), enrichie par l’intelligence artificielle à travers un modèle de langage (LLM) pour la sélection automatique des scanners de sécurité les plus adaptés à chaque projet, ainsi que pour la priorisation contextuelle des vulnérabilités détectées en leur attribuant des scores afin de faciliter leur ordonnancement selon leur criticité. La plateforme intègre également un système de recommandation basé sur le paradigme RAG (Retrieval-Augmented Generation) permettant de proposer des solutions pertinentes pour la correction des vulnérabilités identifiées, en s'appuyant sur le serveur local \textbf{Ollama}. Enfin, elle s’intègre avec des solutions éprouvées de gestion des vulnérabilités, telles que \textbf{DefectDojo} \cite{defectdojo}, afin de centraliser et faciliter le suivi des rapports de sécurité.

Ce rapport est organisé en quatre chapitres:\\
Le \textbf{Chapitre~1} présente le
contexte général du projet, l'organisme d'accueil, la problématique identifiée
ainsi que la solution proposée.\\ Le \textbf{Chapitre~2} détaille la spécification
des besoins à travers les acteurs du système, les besoins fonctionnels et non
fonctionnels, ainsi que les diagrammes de cas d'utilisation.\\ Le \textbf{Chapitre~3}
expose la conception de la solution, en décrivant l'architecture logique, les
diagrammes de classes, le modèle de données et la vue dynamique du système.\\
Enfin, le \textbf{Chapitre~4} présente la phase de réalisation, incluant les
technologies utilisées, les interfaces de l'application et les résultats obtenus.
Une conclusion générale clôture ce rapport en dressant un bilan du travail accompli
et en ouvrant des perspectives d'amélioration.

%============================================================
% CHAPITRE 1 : PRÉSENTATION GÉNÉRALE DU PROJET
%============================================================
\chapter{Présentation générale du projet}
\label{chap:presentation}

\section*{Introduction}
\addcontentsline{toc}{section}{Introduction}

Ce chapitre a pour but de présenter le cadre général dans lequel s'inscrit notre
projet de fin d'études. Nous commençons par décrire l'organisme d'accueil qui nous
a accueillis pour la réalisation de ce stage, puis nous présentons le projet dans
son ensemble : son contexte, la problématique qu'il vise à résoudre, la solution
que nous avons proposée, ainsi que ses objectifs. Nous analysons ensuite les
solutions existantes et nous terminons par une présentation du processus de
développement adopté et de la chronologie du projet.

%------------------------------------------------------------
\section{Présentation de l'organisme d'accueil}
\label{sec:organisme}

Le stage a été réalisé au sein de la société \textbf{Mobelite}, une entreprise spécialisée dans le domaine des nouvelles technologies. 

\subsection{Identité de l'entreprise}
\textbf{Mobelite} est une agence digitale spécialisée dans le conseil en stratégie mobile, la conception et le développement d'applications mobiles, de sites web et le marketing mobile. Forte d'une équipe d'experts, l'entreprise accompagne ses clients dans la réalisation de solutions numériques sur les plateformes les plus répandues (iOS, Android, Web).

\subsection{Fiche de renseignements généraux}
Les principales informations concernant la société sont résumées ci-dessous :
\begin{itemize}
    \item \textbf{Dénomination :} Mobelite (Tunisie / Labs).
    \item \textbf{Directeur Général :} Mohamed Letaief.
    \item \textbf{Sièges sociaux :} Monastir, Sfax et Tunis (Tunisie), Paris (France).
    \item \textbf{Site web :} \url{http://www.mobelite.fr}
    \item \textbf{Secteurs d'activité :} Développement mobile, développement web, Marketing digital.
\end{itemize}

\subsection{Pôles d'expertise et services}
La structure de Mobelite est organisée autour de trois pôles majeurs de compétences :
\begin{itemize}
    \item \textbf{Le Pôle Mobile :} Spécialisé dans l'innovation et le développement d'applications natives et hybrides (iOS, Android, Java/Kotlin) en mettant l'accent sur l'expérience utilisateur et l'optimisation mobile.
    \item \textbf{Le Pôle Web :} Dédié au développement de sites web dynamiques, de plateformes e-commerce et d'applications web modernes intégrant les dernières technologies (HTML5, PHP, architectures réactives).
    \item \textbf{Le Centre de Test :} Une équipe réactive dont la mission est de tester et de valider les projets pour garantir une qualité optimale et un fonctionnement sans faille avant la livraison aux clients.
\end{itemize}

En plus de ses services de conseil et de développement, Mobelite édite ses propres solutions innovantes telles que \textbf{Mobelite Plus} (solution de marketing mobile) et \textbf{Mobelite Publishing} (solution de publication d'app-books). 

Dans le cadre de ses activités, Mobelite accorde une importance capitale à la qualité et à la sécurité des produits logiciels qu'elle développe. Notre stage s'inscrit directement dans cette démarche d'amélioration continue des processus de développement sécurisé (DevSecOps) au sein de l'entreprise.

%------------------------------------------------------------
\section{Présentation du projet}
\label{sec:projet}

\subsection{Cadre général du projet}
\label{subsec:cadre}

Le projet VulnOps s'inscrit dans le domaine de la \textbf{sécurité applicative}
et des \textbf{pratiques DevSecOps}. Il vise à fournir un environnement automatisé
et centralisé pour l'analyse des vulnérabilités dans les projets logiciels hébergés
sur GitHub. La plateforme est destinée à être utilisée par des équipes de développement
souhaitant intégrer la sécurité au sein de leur pipeline de développement logiciel,
sans nécessiter une expertise approfondie en sécurité informatique.

VulnOps fonctionne comme une application web à architecture client-serveur.
Le serveur expose une API REST développée avec \texttt{Django REST Framework},
tandis que le client est une application monopage (\textit{Single Page Application})
développée avec \texttt{React/TypeScript} et construite avec \texttt{Vite}.
La plateforme s'interface avec plusieurs outils d'analyse de sécurité tiers
(Bandit, ESLint, Semgrep, SonarCloud, Detekt, Gosec, etc.)
et utilise le service \textbf{Ollama} pour accéder à des modèles de langage de grande taille (LLM) permettant la mise en place d’un système de recommandation intelligent des outils de scan.
\subsection{Problématique}
\label{subsec:problematique}

Les équipes de développement logiciel font face à plusieurs défis majeurs en
matière de sécurité applicative :

\begin{itemize}
    \item \textbf{La multiplicité des outils :} Il existe de nombreux outils d’analyse de sécurité, chacun étant spécialisé dans un langage ou un type de vulnérabilité particulier. Le choix de l’outil adéquat est souvent complexe et requiert une expertise spécifique. De plus, ces outils sont généralement dispersés et nécessitent une installation et une configuration individuelles, ce qui complique leur adoption. D’où la nécessité d’une plateforme centralisée regroupant ces outils et facilitant leur utilisation sans contrainte d’installation pour l’utilisateur.
    
    \item \textbf{L'intégration fragmentée :} L'utilisation manuelle de ces outils,
    de manière isolée, rend difficile la consolidation et la comparaison des
    résultats, ce qui nuit à une vision globale de la posture de sécurité d'un
    projet.
    
    \item \textbf{Le manque d'automatisation :} Dans de nombreuses organisations,
    les analyses de sécurité sont réalisées de manière ponctuelle et non systématique,
    créant des angles morts qui peuvent être exploités par des acteurs malveillants
    \cite{owasptop10}.
    
    \item \textbf{La gestion des vulnérabilités :} Le suivi des vulnérabilités
    identifiées, de leur priorisation à leur correction, est souvent géré de
    manière informelle, sans traçabilité ni processus établi.
\end{itemize}

\subsection{Solution proposée}
\label{subsec:solution}

En réponse à cette problématique, nous proposons \textbf{VulnOps}, une plateforme
web DevSecOps qui offre les fonctionnalités suivantes :

\begin{enumerate}
    \item \textbf{Analyse multi-scanners automatisée :} VulnOps intègre dix outils
    d'analyse de sécurité (Bandit, ESLint, SonarCloud, Semgrep, Detekt, Gosec,
    Clippy, Psalm, Brakeman, Cppcheck) et permet de les exécuter de manière
    coordonnée sur un dépôt GitHub donné.
    
    \item \textbf{Sélection intelligente des scanners :} En s'appuyant sur un
    modèle de langage (LLM) hébergé localement via \textbf{Ollama}, VulnOps analyse automatiquement
    la structure et les langages d'un projet pour sélectionner les outils de
    scan les mieux adaptés.
    
    \item \textbf{Système de recommandation basé sur RAG :} VulnOps intègre un système de recommandation
    s'appuyant sur le paradigme Retrieval-Augmented Generation (RAG) et un modèle de langage (LLM)
    afin de proposer des solutions pertinentes pour la correction des vulnérabilités détectées,
    ainsi que de prioriser ces dernières en fonction du contexte du projet en leur attribuant
    des scores de criticité.

    \item \textbf{Centralisation et visualisation :} Les résultats de tous les
    scans sont centralisés dans une interface web intuitive, permettant de consulter
    les vulnérabilités détectées, leur sévérité et leur localisation dans le code.
    
   \item \textbf{Intégration avec DefectDojo :} Les rapports de vulnérabilités
    sont automatiquement transmis à DefectDojo \cite{defectdojo}, une plateforme
    open-source de gestion des vulnérabilités, afin d’assurer leur centralisation,
    leur déduplication et leur suivi tout au long de leur cycle de vie. 
    
    \item \textbf{Authentification sécurisée via GitHub OAuth :} L'accès à la
    plateforme est sécurisé par le protocole OAuth~2.0 de GitHub, permettant aux
    utilisateurs de se connecter directement avec leur compte GitHub et d'accéder
    à leurs dépôts.
\end{enumerate}

\subsection{Objectifs}
\label{subsec:objectifs}

Les objectifs du projet VulnOps sont les suivants :

\textbf{Objectifs fonctionnels :}
\begin{itemize}
    \item Permettre l'authentification des utilisateurs via GitHub OAuth~2.0.
    \item Lister les dépôts GitHub de l'utilisateur authentifié.
    \item Permettre le déclenchement et le suivi des analyses de sécurité sur
    un dépôt sélectionné.
    \item Afficher les vulnérabilités détectées avec leur sévérité, localisation
    et description.
    \item Implémenter un mécanisme de sélection automatique des scanners.
        \item Intégrer la plateforme DefectDojo pour la gestion centralisée des
    vulnérabilités.
    \item Implémenter un mécanisme de recommandation basé sur un système RAG (Retrieval-Augmented Generation) et un modèle de langage (LLM) afin de proposer des solutions de correction adaptées aux vulnérabilités détectées.
\end{itemize}

\textbf{Objectifs non fonctionnels :}
\begin{itemize}
    \item \textbf{Performance :} Les analyses doivent s'exécuter en moins de
    dix minutes pour les projets de taille raisonnable.
    \item \textbf{Fiabilité :} La plateforme doit gérer les erreurs de manière
    robuste, avec des mécanismes de repli (\textit{fallback}) appropriés.
    \item \textbf{Sécurité :} Toutes les données sensibles (clés API, tokens)
    doivent être gérées via des variables d'environnement et ne jamais être
    exposées.
    \item \textbf{Maintenabilité :} Le code doit être modulaire, bien documenté
    et facilement extensible pour l'ajout de nouveaux scanners.
    \item \textbf{Ergonomie :} L'interface utilisateur doit être intuitive et
    accessible, même pour des non-spécialistes en sécurité.
\end{itemize}

%------------------------------------------------------------
\section{Étude de l'existant}
\label{sec:existant}

Plusieurs solutions existent sur le marché pour répondre aux besoins d'analyse
de sécurité des applications. Le tableau~\ref{tab:existant} présente une comparaison
de ces solutions avec leurs avantages et inconvénients respectifs.

\begin{table}[H]
    \centering
    \caption{Comparaison des solutions existantes}
    \label{tab:existant}
    \begin{tabularx}{\textwidth}{|>{\bfseries}p{3cm}|X|X|}
        \hline
        \textbf{Système existant} & \textbf{Avantages} & \textbf{Inconvénients} \\
        \hline
        SonarCloud / SonarQube & 
        Analyse multi-langages de qualité ; intégration CI/CD ; tableau de bord riche ; 
        historique des métriques &
        Solution payante pour les fonctionnalités avancées ; nécessite une configuration 
        manuelle par projet ; pas de gestion unifiée des vulnérabilités \\
        \hline
        Snyk & 
        Détection des vulnérabilités dans les dépendances (SCA) ; intégration avec GitHub/GitLab ;
        interface intuitive &
        Fonctionnalités complètes limitées en version gratuite ; pas de SAST intégré
        pour tous les langages ; dépendance à un service tiers payant \\
        \hline
        GitHub Advanced Security (GHAS) &
        Intégration native à GitHub ; analyse des secrets et du code (CodeQL) ;
        alertes Dependabot &
        Réservé aux dépôts privés payants (GitHub Enterprise) ; limité à l'écosystème 
        GitHub ; pas de gestion des vulnérabilités externe \\
        \hline
        OWASP DefectDojo &
        Gestion open-source des vulnérabilités ; supporte de nombreux formats d'outils ;
        suivi du cycle de vie des failles &
        Interface complexe ; ne réalise pas les analyses lui-même ; nécessite une 
        orchestration externe des outils de scan \\
        \hline
        Scripts CI/CD manuels &
        Entièrement personnalisables ; aucun coût de licence &
        Maintenance complexe ; pas d'interface utilisateur dédiée ;
        résultats fragmentés ; expertise technique requise \\
        \hline
    \end{tabularx}
\end{table}

Dans le tableau~\ref{tab:existant}, nous présentons les principales solutions
existantes dans le domaine de l'analyse de sécurité applicative. Cette analyse
met en évidence que, si chaque solution offre des avantages certains, aucune
ne répond de manière complète et abordable à l'ensemble des besoins identifiés :
l'automatisation multi-outils, la sélection intelligente des scanners, la visualisation
unifiée, et la gestion du cycle de vie des vulnérabilités. VulnOps vise précisément
à combler ces lacunes en intégrant ces différentes dimensions en une seule plateforme
cohérente.

%------------------------------------------------------------
\section{Chronologie}
\label{sec:chronologie}

La réalisation du projet VulnOps s'est déroulée sur une période de
\textbf{quatre mois} selon le calendrier présenté dans le
diagramme de Gantt de la figure~\ref{fig:gantt}.

\begin{figure}[H]
    \centering
    \begin{ganttchart}[
        hgrid,
        vgrid,
        x unit=0.65cm,
        y unit title=0.6cm,
        y unit chart=0.6cm,
        title/.style={fill=primaryblue, draw=primaryblue},
        title label font=\small\bfseries\color{white},
        bar/.style={fill=secondaryblue},
        bar label font=\tiny\color{black},
        group/.style={fill=primaryblue},
        group label font=\small\bfseries,
        milestone/.style={fill=primaryblue},
        canvas/.style={draw=black!50, dashed},
    ]{1}{16}
        \gantttitle{Semaines du stage}{16} \\
        \gantttitlelist{1,...,16}{1} \\
        \ganttgroup{Phase 1 : Analyse}{1}{3} \\
        \ganttbar{Étude de l'existant}{1}{2} \\
        \ganttbar{Spécification des besoins}{2}{3} \\
        \ganttgroup{Phase 2 : Conception}{4}{6} \\
        \ganttbar{Architecture logique}{4}{5} \\
        \ganttbar{Modélisation UML}{5}{6} \\
        \ganttbar{Conception BDD}{6}{6} \\
        \ganttgroup{Phase 3 : Développement}{7}{13} \\
        \ganttbar{Backend Django \& API REST}{7}{9} \\
        \ganttbar{Frontend React/TypeScript}{9}{11} \\
        \ganttbar{Intégration scanners}{10}{12} \\
        \ganttbar{Sélection automatique LLM}{12}{13} \\
        \ganttbar{Intégration DefectDojo}{13}{13} \\
        \ganttgroup{Phase 4 : Tests \& validation}{14}{15} \\
        \ganttbar{Tests unitaires \& intégration}{14}{14} \\
        \ganttbar{Correction des bugs}{14}{15} \\
        \ganttgroup{Phase 5 : Rédaction}{15}{16} \\
        \ganttbar{Rapport PFE}{15}{16} \\
    \end{ganttchart}
    \caption{Diagramme de Gantt -- Chronologie du projet VulnOps}
    \label{fig:gantt}
\end{figure}

Dans la figure~\ref{fig:gantt}, nous présentons le diagramme de Gantt illustrant
la planification temporelle du projet VulnOps. Le stage est divisé en cinq phases
principales : l'analyse des besoins, la conception du système, le développement
des différentes composantes, la phase de tests et de validation, et enfin la
rédaction du présent rapport.

%------------------------------------------------------------
\section{Processus de développement}
\label{sec:processus}

Pour la réalisation du projet VulnOps, nous avons adopté une méthodologie de
développement \textbf{agile inspirée de Scrum} \cite{scrumguide}. Cette approche
itérative et incrémentale présente plusieurs avantages dans le cadre de ce projet :

\begin{itemize}
    \item \textbf{Flexibilité face aux changements :} Les exigences ont évolué
    au cours du projet, notamment lors de l'ajout de nouvelles fonctionnalités
    (intégration DefectDojo, module RAG). Scrum permet d'intégrer ces changements
    sans remettre en cause l'ensemble du plan de développement.
    
    \item \textbf{Livraisons incrémentales :} Chaque sprint a abouti à une version
    fonctionnelle et testable de la plateforme, permettant à l'encadrant de valider
    les avancées de manière régulière.
    
    \item \textbf{Gestion des risques :} Les itérations courtes ont permis
    d'identifier et de corriger rapidement les problèmes techniques rencontrés,
    notamment lors de l'intégration des outils de scan.
\end{itemize}

Pour mener à bien ce projet sur une période de \textbf{16 semaines} (quatre mois), le cycle de développement a été organisé en \textbf{phases itératives et incrémentales} calées sur le calendrier du projet. Chaque étape clé (Analyse, Conception, Développement, Tests) a été planifiée pour permettre une validation progressive des avancées de manière régulière. Les fonctionnalités ont été priorisées dans un \textit{product backlog} et assignées à ces différentes phases selon leur importance et leur complexité.

\section*{Conclusion}
\addcontentsline{toc}{section}{Conclusion}

Dans ce chapitre, nous avons présenté le contexte général du projet VulnOps.
Nous avons décrit l'organisme d'accueil, exposé la problématique liée à la 
fragmentation et au manque d'automatisation des outils d'analyse de sécurité, 
et présenté la solution que nous proposons pour y répondre. L'étude de l'existant
a mis en évidence les limites des solutions actuelles et justifié la nécessité 
de développer VulnOps. Le chapitre suivant se consacrera à la spécification 
détaillée des besoins fonctionnels et non fonctionnels de notre système.

%============================================================
% CHAPITRE 2 : SPÉCIFICATION DES BESOINS
%============================================================
\chapter{Spécification des besoins}
\label{chap:besoins}

\section*{Introduction}
\addcontentsline{toc}{section}{Introduction}

Ce chapitre est consacré à la spécification détaillée des besoins du système
VulnOps. Nous commençons par identifier les différents acteurs qui interagissent
avec la plateforme, puis nous décrivons les besoins fonctionnels et non fonctionnels.
Enfin, nous modélisons ces besoins à travers des diagrammes de cas d'utilisation
conformément à la notation UML \cite{umlspec}.

%------------------------------------------------------------
\section{Acteurs}
\label{sec:acteurs}

Le système VulnOps implique deux types d'acteurs principaux :

\begin{enumerate}
    \item \textbf{L'utilisateur authentifié (Développeur / Ingénieur DevSecOps) :}
    Il s'agit de l'acteur principal du système. Après authentification via GitHub
    OAuth~2.0, cet acteur peut accéder à ses dépôts GitHub, déclencher des analyses
    de sécurité, consulter les résultats, interagir avec l'assistant IA et gérer
    ses projets dans la plateforme.
    
    \item \textbf{Le système (acteur secondaire / externe) :}
    Cet acteur représente l'ensemble des services externes avec lesquels VulnOps
    interagit de manière automatisée :
    \begin{itemize}
        \item \textbf{GitHub API :} pour l'authentification OAuth et la récupération
        des dépôts.
        \item \textbf{Les scanners de sécurité :} Bandit, ESLint, Semgrep, SonarCloud,
        Detekt, Gosec, Clippy, Psalm, Brakeman, Cppcheck.
        \item \textbf{Ollama :} pour la sélection intelligente des scanners
        via LLM et le module RAG local.
    \end{itemize}
\end{enumerate}

%------------------------------------------------------------
\section{Besoins fonctionnels}
\label{sec:besoins_fonctionnels}

Les besoins fonctionnels de VulnOps sont les suivants :

\begin{itemize}
    \item \textbf{BF1 -- Authentification GitHub :} L'utilisateur doit pouvoir
    se connecter à la plateforme via son compte GitHub en utilisant le protocole
    OAuth~2.0. Un mode de démonstration sans authentification doit également être
    disponible.
    
    \item \textbf{BF2 -- Gestion des dépôts :} L'utilisateur doit pouvoir
    visualiser la liste de ses dépôts GitHub et sélectionner celui qu'il souhaite
    analyser.
    
    \item \textbf{BF3 -- Déclenchement des analyses :} L'utilisateur doit pouvoir
    sélectionner un ou plusieurs scanners de sécurité et lancer une analyse sur
    le dépôt choisi.
    
    \item \textbf{BF4 -- Sélection automatique des scanners :} Le système doit
    être capable d'analyser automatiquement la structure et les langages d'un
    dépôt afin de recommander les scanners les plus appropriés, en s'appuyant
    sur un modèle de langage (LLM) via Ollama.
    
    \item \textbf{BF5 -- Visualisation des résultats :} L'utilisateur doit pouvoir
    consulter les résultats des analyses, incluant les vulnérabilités détectées,
    leur sévérité (Critique, Haute, Moyenne, Faible, Informationnelle), leur
    localisation dans le code source et leur description.
    
    \item \textbf{BF6 -- Historique des analyses :} L'utilisateur doit pouvoir
    consulter l'historique des analyses précédemment réalisées sur un dépôt.
    
    \item \textbf{BF7 -- Intégration DefectDojo :} Le système doit automatiquement
    transmettre les rapports des analyses terminées à DefectDojo pour la gestion
    centralisée des vulnérabilités.
    
    \item \textbf{BF8 -- Assistant IA (RAG) :} L'utilisateur doit pouvoir
    interagir avec un assistant conversationnel basé sur la technologie RAG pour
    obtenir des explications et des recommandations relatives aux vulnérabilités
    détectées.
\end{itemize}

%------------------------------------------------------------
\section{Besoins non fonctionnels}
\label{sec:besoins_non_fonctionnels}

\begin{itemize}
    \item \textbf{Fiabilité :} Le système doit implémenter des mécanismes de
    repli (\textit{fallback}) pour les services (Ollama,
    DefectDojo). En cas d'indisponibilité d'un service, la plateforme doit
    continuer à fonctionner avec des fonctionnalités dégradées mais opérationnelles.
    
    \item \textbf{Performance :} L'exécution d'une analyse complète sur un projet
    de taille moyenne doit se terminer en moins de dix minutes. Les appels API
    vers les services tiers doivent être soumis à des délais d'expiration de
    trente secondes.
    
    \item \textbf{Sécurité :} Les données sensibles (clés API, tokens GitHub,
    credentials DefectDojo) ne doivent jamais être exposées dans le code source
    ni dans les réponses API. Elles doivent être gérées exclusivement via des
    variables d'environnement. Les en-têtes CORS doivent être configurés
    strictement.
    
    \item \textbf{Maintenabilité :} L'architecture modulaire du backend doit
    permettre l'ajout de nouveaux scanners sans modification des composants
    existants. Chaque scanner dispose de son propre module Python isolé.
    
    \item \textbf{Ergonomie :} L'interface utilisateur doit être accessible à
    des profils non experts en sécurité, avec une navigation claire et des
    visualisations synthétiques des résultats d'analyse.
    
    \item \textbf{Portabilité :} La plateforme doit être déployable sur tout
    système d'exploitation supportant Docker, et les dépendances doivent être
    déclarées de manière explicite et versionnée.
\end{itemize}

%------------------------------------------------------------
\section{Diagrammes de cas d'utilisation}
\label{sec:use_cases}

\subsection{Diagramme général}
\label{subsec:uc_general}

La figure~\ref{fig:uc_general} présente le diagramme de cas d'utilisation général
du système VulnOps. Nous y identifions les trois fonctionnalités principales de
la plateforme : la gestion de l'authentification et des dépôts (CU1), la gestion
des analyses de sécurité (CU2), et la gestion des résultats et du reporting (CU3).

\begin{figure}[H]
    \centering
    \begin{tikzpicture}[
        actor/.style={draw, fill=blue!5, rectangle, rounded corners, minimum width=2cm, minimum height=0.8cm, font=\small},
        usecase/.style={draw, fill=green!5, ellipse, minimum width=3cm, minimum height=1cm, font=\small},
        system/.style={draw, thick, dashed, rectangle, minimum width=10cm, minimum height=7cm},
    ]
        % Système
        \draw[dashed, thick, rounded corners] (-1,-4) rectangle (9.5, 4);
        \node[above, font=\bfseries] at (4.25, 4) {Système VulnOps};

        % Acteur
        \node[actor] (user) at (-3, 0) {Utilisateur};

        % Cas d'utilisation
        \node[usecase] (cu1) at (3, 2.5) {CU1 : Authentification\\et gestion des dépôts};
        \node[usecase] (cu2) at (3, 0) {CU2 : Analyse de\\sécurité multi-scanners};
        \node[usecase] (cu3) at (3, -2.5) {CU3 : Résultats et\\reporting};

        % Acteurs secondaires
        \node[actor] (github) at (8, 2.5) {GitHub API};
        \node[actor] (scanners) at (8, 0) {Scanners \&\\Ollama};
        \node[actor] (defectdojo) at (8, -2.5) {DefectDojo};

        % Flèches acteur → CUs
        \draw[->] (user) -- (cu1);
        \draw[->] (user) -- (cu2);
        \draw[->] (user) -- (cu3);

        % Flèches CUs → acteurs secondaires
        \draw[->] (cu1) -- (github);
        \draw[->] (cu2) -- (scanners);
        \draw[->] (cu3) -- (defectdojo);

    \end{tikzpicture}
    \caption{Diagramme de cas d'utilisation général de VulnOps}
    \label{fig:uc_general}
\end{figure}

Dans la figure~\ref{fig:uc_general}, nous présentons les trois cas d'utilisation
principaux du système. Le cas CU1 regroupe les fonctionnalités d'authentification
et de navigation dans les dépôts, le cas CU2 concerne le cycle complet d'une
analyse de sécurité, et le cas CU3 englobe la consultation des résultats, la
génération de rapports et la transmission à DefectDojo.

\subsection{Diagramme CU1 détaillé -- Authentification et gestion des dépôts}
\label{subsec:uc1}

La figure~\ref{fig:uc1} présente le diagramme détaillé du cas d'utilisation CU1.
Ce cas d'utilisation décrit les interactions de l'utilisateur avec le système
lors de l'authentification et de la consultation de ses dépôts GitHub.

% [Insérer ici le diagramme UML CU1 ou une figure importée]
\begin{figure}[H]
    \centering
    \fbox{\parbox{12cm}{\centering\vspace{2cm}
    \textit{[Insérer ici le diagramme UML de cas d'utilisation CU1 --}\\
    \textit{Authentification GitHub OAuth et gestion des dépôts]}
    \vspace{2cm}}}
    \caption{Diagramme de cas d'utilisation CU1 -- Authentification et gestion des dépôts}
    \label{fig:uc1}
\end{figure}

\textbf{Description textuelle du cas d'utilisation principal : Se connecter via GitHub}

\begin{table}[H]
    \centering
    \caption{Description textuelle -- CU1 : Se connecter via GitHub OAuth}
    \label{tab:desc_uc1}
    \begin{tabularx}{\textwidth}{|l|X|}
        \hline
        \textbf{Identifiant} & CU1-01 \\
        \hline
        \textbf{Nom} & Se connecter via GitHub OAuth \\
        \hline
        \textbf{Acteur principal} & Utilisateur \\
        \hline
        \textbf{Préconditions} & L'utilisateur possède un compte GitHub valide \\
        \hline
        \textbf{Postconditions} & L'utilisateur est authentifié et accède au tableau de bord \\
        \hline
        \textbf{Scénario nominal} & 
        \begin{enumerate}[leftmargin=1.2em]
            \item L'utilisateur accède à la page de connexion de VulnOps.
            \item L'utilisateur clique sur "Se connecter avec GitHub".
            \item VulnOps redirige l'utilisateur vers la page d'autorisation GitHub.
            \item L'utilisateur autorise l'application VulnOps sur GitHub.
            \item GitHub retourne un code d'autorisation à VulnOps.
            \item Le backend VulnOps échange ce code contre un token d'accès.
            \item Le système crée ou met à jour le profil utilisateur en base de données.
            \item L'utilisateur est redirigé vers le tableau de bord de VulnOps.
        \end{enumerate} \\
        \hline
        \textbf{Scénarios alternatifs} & 
        L'utilisateur refuse l'autorisation GitHub : le système affiche un message d'erreur et propose une connexion en mode démonstration. \\
        \hline
    \end{tabularx}
\end{table}

Dans le tableau~\ref{tab:desc_uc1}, nous présentons la description textuelle
du scénario nominal d'authentification via GitHub OAuth. Ce processus suit le
flux standard du protocole OAuth~2.0, avec un échange sécurisé du code temporaire
contre un token d'accès permanent \cite{oauth2rfc}.

\subsection{Diagramme CU2 détaillé -- Analyse de sécurité multi-scanners}
\label{subsec:uc2}

Le cas d'utilisation CU2 décrit le processus complet d'une analyse de sécurité
sur un dépôt GitHub. La figure~\ref{fig:uc2} illustre les différentes actions
possibles pour l'utilisateur dans ce contexte.

\begin{figure}[H]
    \centering
    \fbox{\parbox{12cm}{\centering\vspace{2cm}
    \textit{[Insérer ici le diagramme UML de cas d'utilisation CU2 --}\\
    \textit{Sélection du scanner, déclenchement et suivi du scan]}
    \vspace{2cm}}}
    \caption{Diagramme de cas d'utilisation CU2 -- Analyse de sécurité multi-scanners}
    \label{fig:uc2}
\end{figure}

\textbf{Diagramme de séquence système -- Lancer une analyse automatique}

La figure~\ref{fig:seq_autoscan} présente le diagramme de séquence système
pour le cas nominal du déclenchement d'une analyse automatique avec sélection
intelligente des scanners.

\begin{figure}[H]
    \centering
    \fbox{\parbox{12cm}{\centering\vspace{2cm}
    \textit{[Insérer ici le diagramme de séquence système --}\\
    \textit{Auto-sélection des scanners et lancement du scan]}
    \vspace{2cm}}}
    \caption{Diagramme de séquence système -- Lancement d'une analyse automatique}
    \label{fig:seq_autoscan}
\end{figure}

Dans la figure~\ref{fig:seq_autoscan}, nous illustrons la séquence d'interactions
entre l'utilisateur, le frontend React, le backend Django et les services externes
(GitHub, Ollama, Scanners) lors du déclenchement d'une analyse automatique.

\subsection{Diagramme CU3 détaillé -- Résultats et reporting}
\label{subsec:uc3}

Le cas CU3 couvre la consultation des résultats d'analyse, la génération de
rapports consolidés et la transmission des données à DefectDojo.

\begin{figure}[H]
    \centering
    \fbox{\parbox{12cm}{\centering\vspace{2cm}
    \textit{[Insérer ici le diagramme UML de cas d'utilisation CU3 --}\\
    \textit{Consultation des résultats, rapport final et IntégrationDefectDojo]}
    \vspace{2cm}}}
    \caption{Diagramme de cas d'utilisation CU3 -- Résultats et reporting}
    \label{fig:uc3}
\end{figure}

\section*{Conclusion}
\addcontentsline{toc}{section}{Conclusion}

Ce chapitre a permis d'identifier les acteurs du système VulnOps et de spécifier
de manière détaillée les besoins fonctionnels et non fonctionnels de la plateforme.
La modélisation UML à travers les diagrammes de cas d'utilisation a mis en évidence
les trois grandes fonctionnalités du système : la gestion de l'authentification
et des dépôts, la réalisation des analyses de sécurité multi-scanners, et
la consultation des résultats avec leur intégration dans DefectDojo.
Le chapitre suivant sera consacré à la conception de la solution, en détaillant
l'architecture logique du système, les diagrammes de classes et le modèle
de données.

%============================================================
% CHAPITRE 3 : CONCEPTION
%============================================================
\chapter{Conception}
\label{chap:conception}

\section*{Introduction}
\addcontentsline{toc}{section}{Introduction}

Ce chapitre présente la phase de conception de la plateforme VulnOps. Nous
décrivons successivement la vue statique de l'application, à travers l'architecture
logique, les diagrammes de classes et le modèle de données, puis la vue dynamique,
à travers des diagrammes de séquence, d'activités et d'états-transitions.
L'objectif est de fournir une représentation structurée et formelle du système
avant son implémentation.

%------------------------------------------------------------
\section{Vue statique de l'application}

\subsection{Architecture logique du logiciel}
\label{subsec:architecture_logique}

L'architecture logique de VulnOps s'appuie sur le patron architectural
\textbf{Clean Architecture} \cite{cleanarchitecture}, qui favorise la séparation
des préoccupations et l'indépendance des couches applicatives. Par ailleurs,
le backend Django suit naturellement le patron \textbf{MVT} (Model-View-Template),
adapté ici à une API REST en \textbf{MVC} (Model-View-Controller).

La figure~\ref{fig:architecture} présente le diagramme de paquetages représentant
l'architecture logique du système VulnOps.

\begin{figure}[H]
    \centering
    \begin{tikzpicture}[
        package/.style={draw, rectangle, rounded corners=2pt, minimum width=3.5cm,
                        minimum height=1.2cm, fill=blue!8, font=\small\bfseries},
        arrow/.style={->, thick, >=stealth},
    ]
        % Frontend
        \node[package, fill=green!10] (ui) at (0, 8) {Couche Présentation\\(React / TypeScript)};

        % API Gateway
        \node[package, fill=yellow!15] (api) at (0, 5.5) {Couche API\\(Django REST Framework)};

        % Services
        \node[package, fill=blue!10] (auth) at (-5, 3) {Module\\Authentification\\(GitHub OAuth)};
        \node[package, fill=blue!10] (scanner) at (0, 3) {Module\\Scanner\\(Orchestrateur)};
        \node[package, fill=blue!10] (rag) at (5, 3) {Module RAG\\(Ollama / LLM)};

        % Sous-modules scanner
        \node[package, fill=orange!10] (bandit) at (-3.5, 0.5) {Bandit};
        \node[package, fill=orange!10] (eslint) at (-1, 0.5) {ESLint};
        \node[package, fill=orange!10] (semgrep) at (1.5, 0.5) {Semgrep};
        \node[package, fill=orange!10] (autres) at (4, 0.5) {Autres\\scanners...};

        % Couche données
        \node[package, fill=red!10] (db) at (-3, -2) {Base de données\\(SQLite / PostgreSQL)};
        \node[package, fill=red!10] (defectdojo) at (3, -2) {DefectDojo\\(externe)};

        % Ollama
        \node[package, fill=purple!10] (ollama) at (0, -2) {Ollama\\(LLM sélection)};

        % Flèches
        \draw[arrow] (ui) -- (api);
        \draw[arrow] (api) -- (auth);
        \draw[arrow] (api) -- (scanner);
        \draw[arrow] (api) -- (rag);
        \draw[arrow] (scanner) -- (bandit);
        \draw[arrow] (scanner) -- (eslint);
        \draw[arrow] (scanner) -- (semgrep);
        \draw[arrow] (scanner) -- (autres);
        \draw[arrow] (scanner) -- (ollama);
        \draw[arrow] (api) -- (db);
        \draw[arrow] (api) -- (defectdojo);
    \end{tikzpicture}
    \caption{Diagramme de paquetages -- Architecture logique de VulnOps}
    \label{fig:architecture}
\end{figure}

Dans la figure~\ref{fig:architecture}, nous présentons l'architecture logique
de VulnOps organisée en couches distinctes. La couche de présentation, développée
en React/TypeScript, communique exclusivement avec la couche API via des requêtes
HTTP REST. Le backend Django orchestre trois modules principaux : le module
d'authentification via GitHub OAuth, le module de scanning qui encapsule les
différents outils d'analyse, et le module RAG pour l'assistant conversationnel.
La couche de persistance utilise SQLite en développement et peut être migrée
vers PostgreSQL en production.

\subsection{Diagrammes de classes}
\label{subsec:classes}

La figure~\ref{fig:classes_scanner} présente le diagramme de classes du module
Scanner, qui constitue le cœur fonctionnel de la plateforme VulnOps.

\begin{figure}[H]
    \centering
    \fbox{\parbox{14cm}{\centering\vspace{3cm}
    \textit{[Insérer ici le diagramme de classes UML du module Scanner --}\\
    \textit{Classes : ScanResult, Vulnerability, ProjectAnalyzer,}\\
    \textit{AutoScannerOrchestrator, OllamaSelector, et les runners]}
    \vspace{3cm}}}
    \caption{Diagramme de classes -- Module Scanner de VulnOps}
    \label{fig:classes_scanner}
\end{figure}

Dans la figure~\ref{fig:classes_scanner}, nous présentons le diagramme de classes
du module Scanner. Les principales classes sont les suivantes :

\begin{itemize}
    \item \textbf{ScanResult :} Modèle Django représentant le résultat d'une
    analyse. Ses attributs principaux sont \texttt{scan\_id}, \texttt{repo\_name},
    \texttt{scanner\_type}, \texttt{status}, \texttt{created\_at},
    \texttt{total\_issues}, \texttt{high\_count}, \texttt{medium\_count},
    \texttt{low\_count}. Cette classe expose les méthodes \texttt{get\_vulnerabilities()}
    et \texttt{get\_metrics()}.
    
    \item \textbf{Vulnerability :} Représente une vulnérabilité individuelle
    détectée lors d'une analyse. Ses attributs incluent \texttt{test\_id},
    \texttt{severity}, \texttt{confidence}, \texttt{filename},
    \texttt{line\_number}, \texttt{issue\_text} et \texttt{cwe}.
    
    \item \textbf{ProjectAnalyzer :} Classe responsable de l'analyse de la
    structure d'un projet cloné. Ses méthodes principales sont
    \texttt{analyze()} (retourne un dictionnaire des langages et frameworks
    détectés) et \texttt{get\_scan\_candidates()}.
    
    \item \textbf{OllamaSelector :} Classe implémentant le patron de
    conception \textit{Strategy} \cite{gof} pour la sélection des scanners.
    Elle expose la méthode \texttt{select\_scanners(analysis)} qui appelle
    le serveur local Ollama et retourne une liste de scanners recommandés avec leur
    score de confiance.
    
    \item \textbf{AutoScannerOrchestrator :} Classe coordinatrice qui orchestre
    le flux complet : clonage du dépôt, analyse, sélection et lancement des
    scanners.
\end{itemize}

\subsection{Bases de données}
\label{subsec:bdd}

La figure~\ref{fig:erd} présente le modèle entités-associations de la base
de données de VulnOps.

\begin{figure}[H]
    \centering
    \fbox{\parbox{14cm}{\centering\vspace{3cm}
    \textit{[Insérer ici le diagramme entités-associations --}\\
    \textit{Entités : User, GitHubProfile, Project, ScanResult, Vulnerability, DefectDojoReport]}
    \vspace{3cm}}}
    \caption{Diagramme entités-associations -- Base de données VulnOps}
    \label{fig:erd}
\end{figure}

Dans la figure~\ref{fig:erd}, nous présentons le modèle de données de VulnOps.
Les entités principales sont :

\begin{itemize}
    \item \textbf{User :} Utilisateur de la plateforme (lié au modèle
    \texttt{AbstractUser} de Django). Attributs : \texttt{id}, \texttt{username},
    \texttt{email}, \texttt{date\_joined}.
    
    \item \textbf{GitHubProfile :} Extension du modèle User pour stocker
    les informations GitHub. Attributs : \texttt{github\_id},
    \texttt{github\_username}, \texttt{github\_token}, \texttt{avatar\_url}.
    
    \item \textbf{ScanResult :} Résultat d'une analyse de sécurité. Attributs :
    \texttt{scan\_id}, \texttt{repo\_name}, \texttt{repo\_full\_name},
    \texttt{scanner\_type}, \texttt{status}, \texttt{created\_at},
    \texttt{total\_issues}, \texttt{vulnerabilities} (JSON).
\end{itemize}

%------------------------------------------------------------
\section{Vue dynamique de l'application}

\subsection{Diagrammes de séquences}
\label{subsec:seq}

La figure~\ref{fig:seq_scan} présente le diagramme de séquence détaillé
du processus d'analyse de sécurité avec sélection automatique des scanners.
Ce scénario constitue la fonctionnalité principale de VulnOps.

\begin{figure}[H]
    \centering
    \fbox{\parbox{14cm}{\centering\vspace{3cm}
    \textit{[Insérer ici le diagramme de séquence UML --}\\
    \textit{Flux : Frontend → API → ProjectAnalyzer → Ollama → Scanner → DefectDojo]}
    \vspace{3cm}}}
    \caption{Diagramme de séquence -- Analyse automatique avec sélection LLM}
    \label{fig:seq_scan}
\end{figure}

Dans la figure~\ref{fig:seq_scan}, nous décrivons la séquence d'interactions
lors d'une analyse automatique. L'utilisateur soumet l'URL d'un dépôt GitHub ;
le backend clone le dépôt, analyse sa structure via \texttt{ProjectAnalyzer},
interroge Ollama pour la sélection intelligente des scanners, exécute
les scanners retenus en parallèle, consolide les résultats, les stocke en base
de données et les transmet à DefectDojo.

\subsection{Diagrammes d'activités}
\label{subsec:activites}

La figure~\ref{fig:activity_scan} illustre le diagramme d'activités du processus
d'analyse de sécurité, depuis la réception de la requête jusqu'à la notification
de fin d'analyse.

\begin{figure}[H]
    \centering
    \fbox{\parbox{14cm}{\centering\vspace{3cm}
    \textit{[Insérer ici le diagramme d'activités UML --}\\
    \textit{Processus de scan : clonage → analyse → sélection →}\\
    \textit{exécution parallèle → agrégation résultats → stockage → notification]}
    \vspace{3cm}}}
    \caption{Diagramme d'activités -- Processus d'analyse de sécurité}
    \label{fig:activity_scan}
\end{figure}

\subsection{Diagrammes états-transitions}
\label{subsec:etats}

La figure~\ref{fig:statechart} présente le diagramme d'états-transitions
de l'entité \texttt{ScanResult}, qui est l'objet central du système VulnOps.

\begin{figure}[H]
    \centering
    \begin{tikzpicture}[
        state/.style={draw, rounded rectangle, fill=blue!10, minimum width=2.5cm,
                      minimum height=0.8cm, font=\small},
        arrow/.style={->, thick, >=stealth},
    ]
        \node[state, fill=black, text=white, minimum width=0.5cm, minimum height=0.5cm] (start) at (0, 0) {};
        \node[state] (pending) at (3, 0) {EN ATTENTE};
        \node[state] (running) at (7, 0) {EN COURS};
        \node[state, fill=green!20] (completed) at (7, -3) {TERMINÉ};
        \node[state, fill=red!20] (failed) at (3, -3) {ÉCHEC};

        \draw[arrow] (start) -- (pending);
        \draw[arrow] (pending) -- node[above, font=\tiny]{déclenchement du scan} (running);
        \draw[arrow] (running) -- node[right, font=\tiny]{succès} (completed);
        \draw[arrow] (running) -- node[above, font=\tiny]{erreur} (failed);
        \draw[arrow] (failed) -- node[below, font=\tiny]{réessayer} (pending);
    \end{tikzpicture}
    \caption{Diagramme états-transitions -- Objet ScanResult}
    \label{fig:statechart}
\end{figure}

Dans la figure~\ref{fig:statechart}, nous présentons les quatre états possibles
d'un objet \texttt{ScanResult} dans VulnOps. Un scan est initialement créé
dans l'état \texttt{EN ATTENTE}, puis passe à l'état \texttt{EN COURS} lorsque
l'exécution est déclenchée. En cas de succès, il atteint l'état \texttt{TERMINÉ}
et les résultats sont disponibles ; en cas d'erreur, il passe à l'état
\texttt{ÉCHEC} avec possibility de relance.

\section*{Conclusion}
\addcontentsline{toc}{section}{Conclusion}

Ce chapitre a présenté la conception complète de la plateforme VulnOps.
L'architecture logique en couches, inspirée de la \textit{Clean Architecture},
assure une séparation claire des responsabilités et facilite la maintenance
et l'extension du système. Les diagrammes de classes ont mis en évidence
les principales entités et leurs relations, tandis que les vues dynamiques
ont illustré les flux de traitement principaux. Le chapitre suivant sera
consacré à la phase de réalisation, présentant les technologies utilisées
et les interfaces de l'application développée.

%============================================================
% CHAPITRE 4 : RÉALISATION
%============================================================
\chapter{Réalisation}
\label{chap:realisation}

\section*{Introduction}
\addcontentsline{toc}{section}{Introduction}

Ce chapitre présente les aspects pratiques de la mise en œuvre de la plateforme
VulnOps. Nous commençons par décrire les technologies et les outils utilisés
pour le développement, puis nous présentons l'architecture physique du système
déployé. Nous détaillons ensuite les principales interfaces de l'application
à travers des captures d'écran illustrées et commentées. Nous terminons par
un bilan des résultats obtenus et les difficultés rencontrées.

%------------------------------------------------------------
\section{Technologies}
\label{sec:technologies}

Le tableau~\ref{tab:technologies} récapitule les principales technologies
utilisées pour le développement de VulnOps.

\begin{table}[H]
    \centering
    \caption{Technologies utilisées dans le développement de VulnOps}
    \label{tab:technologies}
    \begin{tabularx}{\textwidth}{|l|l|X|}
        \hline
        \textbf{Catégorie} & \textbf{Technologie} & \textbf{Rôle / Justification} \\
        \hline
        \multirow{3}{*}{Backend} 
            & Python 3.11 & Langage principal du serveur, riche écosystème sécurité \\
            \cline{2-3}
            & Django 4.2 & Framework web MVT, ORM intégré, gestion des migrations \\
            \cline{2-3}
            & Django REST Framework 3.15 & Construction des endpoints API REST avec sérialisation \\
        \hline
        \multirow{3}{*}{Frontend}
            & React 18 & Framework SPA, composants réutilisables, Virtual DOM \\
            \cline{2-3}
            & TypeScript & Typage statique, meilleure maintenabilité du code \\
            \cline{2-3}
            & Vite & Bundler ultra-rapide, Hot Module Replacement (HMR) \\
        \hline
        \multirow{5}{*}{Scanners SAST}
            & Bandit & Analyse statique Python (vulnérabilités OWASP) \\
            \cline{2-3}
            & ESLint & Analyse statique JavaScript/TypeScript \\
            \cline{2-3}
            & Semgrep 1.45 & Patterns multi-langages, règles OWASP Top 10 \\
            \cline{2-3}
            & SonarCloud & Analyse qualité et sécurité multi-langages en cloud \\
            \cline{2-3}
            & Detekt, Gosec, Clippy, Psalm, Brakeman, Cppcheck & Scanners spécialisés par langage \\
        \hline
        \multirow{2}{*}{IA / LLM}
            & Ollama & Accès aux LLMs locaux (Mistral 7B) pour sélection automatique \\
            \cline{2-3}
            & Ollama & Inférence locale pour le module RAG \\
        \hline
        \multirow{2}{*}{Sécurité}
            & GitHub OAuth 2.0 & Authentification sécurisée via le protocole standard \\
            \cline{2-3}
            & PyGithub 2.4 & SDK Python pour l'API GitHub \\
        \hline
        Vuln. Mgmt & DefectDojo & Plateforme open-source de gestion des vulnérabilités \\
        \hline
        Base de données & SQLite & Développement ; PostgreSQL recommandé en production \\
        \hline
        Versionnement & Git / GitHub & Gestion du code source et collaboration \\
        \hline
    \end{tabularx}
\end{table}

Dans le tableau~\ref{tab:technologies}, nous récapitulons l'ensemble des
technologies mobilisées pour développer VulnOps. Le choix de Django comme
framework backend est justifié par sa maturité, son ORM puissant et sa capacité
à développer rapidement des APIs REST robustes via Django REST Framework.
React a été retenu pour le frontend en raison de sa flexibilité et de la richesse
de son écosystème \cite{reactdocs}.

%------------------------------------------------------------
\section{Outils d'implémentation}
\label{sec:outils}

\begin{table}[H]
    \centering
    \caption{Outils d'implémentation utilisés}
    \label{tab:outils}
    \begin{tabularx}{\textwidth}{|l|X|}
        \hline
        \textbf{Outil} & \textbf{Description} \\
        \hline
        Visual Studio Code & IDE principal pour le développement frontend et backend,
        avec extensions Python, TypeScript et ESLint \\
        \hline
        Postman & Test et validation des endpoints API REST \\
        \hline
        Docker \& Docker Compose & Conteneurisation de DefectDojo et orchestration
        des services locaux \\
        \hline
        Git (v2.44) & Gestion de versions et suivi des modifications \\
        \hline
        Node.js 20 (npm) & Gestion des dépendances JavaScript et build du frontend \\
        \hline
        pip / venv & Gestion des dépendances Python et isolation de l'environnement \\
        \hline
        \textbf{Matériel} & Processeur : [CPU], RAM : [RAM Go], OS : [Windows/Linux] \\
        \hline
    \end{tabularx}
\end{table}

Dans le tableau~\ref{tab:outils}, nous décrivons les outils matériels et
logiciels utilisés lors du développement de VulnOps. Docker a joué un rôle
particulièrement important pour le déploiement local de DefectDojo, permettant
d'isoler ses dépendances complexes du reste de l'environnement de développement.

%------------------------------------------------------------
\section{Architecture physique et évolution temporelle}
\label{sec:arch_physique}

La figure~\ref{fig:deploiement} présente le diagramme de déploiement de
VulnOps, illustrant la superposition de l'architecture logique et des
composants physiques du système.

\begin{figure}[H]
    \centering
    \begin{tikzpicture}[
        node1/.style={draw, fill=blue!10, rectangle, rounded corners,
                      minimum width=4cm, minimum height=1cm, font=\small},
        node2/.style={draw, fill=green!10, rectangle, rounded corners,
                      minimum width=3.5cm, minimum height=0.9cm, font=\small},
        box/.style={draw, thick, dashed, rectangle, rounded corners,
                    fill=gray!5, font=\small\bfseries},
        arrowD/.style={<->, thick, >=stealth, dashed},
    ]
        % Poste développeur / utilisateur
        \draw[thick] (-5, -3.5) rectangle (-0.5, 3.5);
        \node[above, font=\small\bfseries] at (-2.75, 3.5) {Poste Client (Navigateur)};
        \node[node1] at (-2.75, 2.5) {React / TypeScript\\(Vite SPA)};
        \node[node1, fill=yellow!15] at (-2.75, 1) {Axios (HTTP Client)};

        % Serveur Backend
        \draw[thick] (1, -3.5) rectangle (6.5, 3.5);
        \node[above, font=\small\bfseries] at (3.75, 3.5) {Serveur Backend (Python)};
        \node[node1, fill=orange!15] at (3.75, 2.5) {Django REST Framework};
        \node[node1] at (3.75, 1) {Module Scanner\\(Orchestrateur)};
        \node[node1] at (3.75, -0.5) {SQLite / PostgreSQL};
        \node[node2, fill=purple!10] at (3.75, -2) {Bandit | ESLint | Semgrep...};

        % Services externes
        \draw[thick] (8, 0) rectangle (12, 3.5);
        \node[above, font=\small\bfseries] at (10, 3.5) {Services Externes};
        \node[node2] at (10, 2.5) {GitHub API};
        \node[node2, fill=orange!10] at (10, 1) {Ollama (LLM)};

        % DefectDojo (Docker)
        \draw[thick] (8, -3.5) rectangle (12, -0.5);
        \node[above, font=\small\bfseries] at (10, -0.5) {Docker (Local)};
        \node[node2, fill=red!10] at (10, -1.5) {DefectDojo};
        \node[node2] at (10, -2.8) {PostgreSQL (Dojo)};

        % Flèches
        \draw[arrowD] (-0.5, 1.5) -- (1, 1.5) node[midway, above, font=\tiny]{HTTP/REST};
        \draw[arrowD] (6.5, 2.5) -- (8, 2.5) node[midway, above, font=\tiny]{HTTPS};
        \draw[arrowD] (6.5, 1) -- (8, 1) node[midway, above, font=\tiny]{HTTPS};
        \draw[arrowD] (6.5, -1.5) -- (8, -1.5) node[midway, above, font=\tiny]{HTTP};
    \end{tikzpicture}
    \caption{Diagramme de déploiement -- Architecture physique de VulnOps}
    \label{fig:deploiement}
\end{figure}

Dans la figure~\ref{fig:deploiement}, nous présentons l'architecture physique
de VulnOps. Le frontend React s'exécute dans le navigateur de l'utilisateur
et communique avec le backend Django via des requêtes HTTP REST. Le backend
orchestre les différents outils de scan, communique avec les APIs externes
(GitHub, Ollama) et transmet les résultats à DefectDojo déployé localement
via Docker.

%------------------------------------------------------------
\section{Interfaces de l'application}
\label{sec:interfaces}

\subsection{Page d'authentification}

La figure~\ref{fig:ui_login} présente la page d'authentification de VulnOps.

\begin{figure}[H]
    \centering
    \fbox{\parbox{13cm}{\centering\vspace{3cm}
    \textit{[Insérer ici une capture d'écran de la page de connexion --}\\
    \textit{Bouton "Se connecter avec GitHub" et mode démonstration]}
    \vspace{3cm}}}
    \caption{Interface de la page d'authentification de VulnOps}
    \label{fig:ui_login}
\end{figure}

Dans la figure~\ref{fig:ui_login}, nous présentons l'interface de connexion
à la plateforme VulnOps. L'utilisateur dispose de deux options : se connecter
via son compte GitHub grâce au protocole OAuth~2.0, ou accéder à la plateforme
en mode démonstration avec des données fictives.

\subsection{Tableau de bord et liste des dépôts}

La figure~\ref{fig:ui_dashboard} présente le tableau de bord principal affichant
la liste des dépôts GitHub de l'utilisateur authentifié.

\begin{figure}[H]
    \centering
    \fbox{\parbox{13cm}{\centering\vspace{3cm}
    \textit{[Insérer ici une capture d'écran du tableau de bord --}\\
    \textit{Liste des dépôts GitHub avec statistiques sommaires]}
    \vspace{3cm}}}
    \caption{Tableau de bord VulnOps -- Liste des dépôts GitHub}
    \label{fig:ui_dashboard}
\end{figure}

\subsection{Page d'analyse de sécurité}

La figure~\ref{fig:ui_analysis} illustre la page d'analyse d'un dépôt,
permettant la sélection des scanners et le déclenchement des analyses.

\begin{figure}[H]
    \centering
    \fbox{\parbox{13cm}{\centering\vspace{3cm}
    \textit{[Insérer ici une capture d'écran de la page d'analyse --}\\
    \textit{Sélection des scanners, bouton "Auto-Scan", statut en cours]}
    \vspace{3cm}}}
    \caption{Page d'analyse -- Sélection des scanners et déclenchement du scan}
    \label{fig:ui_analysis}
\end{figure}

\subsection{Visualisation des résultats}

La figure~\ref{fig:ui_results} présente l'interface de visualisation des
vulnérabilités détectées lors d'une analyse.

\begin{figure}[H]
    \centering
    \fbox{\parbox{13cm}{\centering\vspace{3cm}
    \textit{[Insérer ici une capture d'écran des résultats d'analyse --}\\
    \textit{Liste des vulnérabilités avec sévérité, localisation et description]}
    \vspace{3cm}}}
    \caption{Interface des résultats d'analyse -- Vulnérabilités détectées}
    \label{fig:ui_results}
\end{figure}

\subsection{Rapport final consolidé (DefectDojo)}

La figure~\ref{fig:ui_rapport} montre l'onglet "Rapport Final" regroupant
les données consolidées transmises à DefectDojo.

\begin{figure}[H]
    \centering
    \fbox{\parbox{13cm}{\centering\vspace{3cm}
    \textit{[Insérer ici une capture d'écran de l'onglet Rapport Final --}\\
    \textit{Statistiques globales, répartition par sévérité, findings DefectDojo]}
    \vspace{3cm}}}
    \caption{Rapport final consolidé -- Intégration DefectDojo}
    \label{fig:ui_rapport}
\end{figure}

%------------------------------------------------------------
\section{Résultats et statistiques}
\label{sec:resultats}

Afin de valider l'efficacité de la plateforme VulnOps, nous avons réalisé
des analyses sur plusieurs dépôts GitHub représentatifs. Le tableau~\ref{tab:resultats}
présente les résultats obtenus lors de ces tests.

\begin{table}[H]
    \centering
    \caption{Résultats des analyses de sécurité sur des projets tests}
    \label{tab:resultats}
    \begin{tabularx}{\textwidth}{|l|l|c|c|c|c|}
        \hline
        \textbf{Projet} & \textbf{Scanner(s)} & \textbf{Critique} & \textbf{Haute} & \textbf{Moyenne} & \textbf{Faible} \\
        \hline
        Projet Python test & Bandit + Semgrep & 0 & 5 & 12 & 8 \\
        \hline
        Projet JS/TS test & ESLint + Semgrep & 0 & 2 & 7 & 15 \\
        \hline
        Projet Java test & SonarCloud & 1 & 3 & 10 & 6 \\
        \hline
        Projet multi-lang & Auto-sélection LLM & 1 & 8 & 18 & 22 \\
        \hline
    \end{tabularx}
\end{table}

Dans le tableau~\ref{tab:resultats}, nous présentons les résultats des analyses
effectuées sur quatre projets de test représentatifs des cas d'usage principaux
de VulnOps. Ces résultats démontrent la capacité de la plateforme à détecter
des vulnérabilités de différents niveaux de sévérité dans des projets de
nature diverse.

Le module de sélection automatique des scanners, basé sur le modèle de langage
Mistral 7B via Ollama, a atteint un taux de pertinence de \textbf{[X\%]}
dans le choix des outils de scan appropriés, validé par comparaison avec
une sélection manuelle effectuée par un expert en sécurité.

%------------------------------------------------------------
\section{Problèmes rencontrés}
\label{sec:problemes}

Au cours du développement de VulnOps, nous avons été confrontés à plusieurs
défis techniques dont voici les principaux :

\begin{enumerate}
    \item \textbf{Configuration CORS :} La communication entre le frontend
    (port~5173) et le backend (port~8000) a nécessité une configuration précise
    des en-têtes CORS dans Django, notamment lors de la gestion des cookies de
    session transmis avec les requêtes \texttt{credentials}. La solution a été
    d'utiliser \texttt{django-cors-headers} avec une configuration explicite
    des origines autorisées et l'activation de \texttt{CORS\_ALLOW\_CREDENTIALS}.
    
    \item \textbf{Hétérogénéité des formats de sortie des scanners :}
    Chaque outil de scan produit ses résultats dans un format JSON spécifique.
    L'implémentation d'un format de sortie uniforme a requis le développement
    de parseurs dédiés pour chacun des dix scanners intégrés, ce qui a représenté
    une charge de travail significative.
    
    \item \textbf{Timeout lors des analyses de grande envergure :}
    Pour les dépôts volumineux, l'exécution de certains scanners (notamment SonarCloud
    et Semgrep) peut dépasser les délais d'expiration configurés. Nous avons
    résolu ce problème en implémentant des \textit{shallow clones} de Git
    (profondeur~1) et en configurant des timeouts spécifiques par scanner.
    
    \item \textbf{Authentification GitHub OAuth en développement local :}
    La configuration de l'OAuth GitHub pour un environnement local a nécessité
    la création d'une application GitHub de développement avec des URLs de callback
    spécifiques, et la gestion sécurisée des tokens via des variables d'environnement.
    
    \item \textbf{Intégration DefectDojo :}
    Le déploiement de DefectDojo via Docker Compose et son initialisation ont
    présenté des complexités dues à sa pile technologique (Django + Celery +
    Redis + PostgreSQL). La configuration de l'API DefectDojo et la gestion
    des engagements et findings ont nécessité une compréhension approfondie
    de son modèle de données.
\end{enumerate}

\section*{Conclusion}
\addcontentsline{toc}{section}{Conclusion}

Ce chapitre a présenté la phase de réalisation de la plateforme VulnOps.
Nous avons détaillé les technologies et outils utilisés, décrit l'architecture
physique de déploiement et illustré les principales interfaces de l'application.
Les résultats obtenus lors des tests ont confirmé la capacité de la plateforme
à détecter des vulnérabilités dans des projets de natures diverses, tandis que le
module de sélection automatique par LLM a démontré son utilité pour guider les
développeurs vers les outils les plus adaptés. Les problèmes rencontrés, bien que
significatifs, ont constitué d'excellentes opportunités d'apprentissage et ont
contribué à renforcer la robustesse de la solution finale.

%============================================================
% CONCLUSION GÉNÉRALE
%============================================================
\chapter*{Conclusion Générale}
\addcontentsline{toc}{chapter}{Conclusion Générale}
\markboth{Conclusion Générale}{Conclusion Générale}

Le travail réalisé dans le cadre de ce projet de fin d'études s'inscrit dans
le domaine de la sécurité applicative et des pratiques DevSecOps. La plateforme
VulnOps a été développée pour répondre au besoin croissant des équipes de
développement logiciel d'intégrer la sécurité de manière automatisée et centralisée
dans leur cycle de développement. L'organisme d'accueil, \textbf{Mobelite},
a fourni un cadre approprié pour l'identification des besoins réels et la
validation des fonctionnalités développées.

Au terme de ce stage, nous avons réalisé une plateforme web fonctionnelle qui
intègre 12 outils d'analyse de sécurité (Bandit, ESLint, Semgrep, SonarCloud,
Detekt, Gosec, Clippy, Psalm, Brakeman, Cppcheck, Trivy et OWASP ZAP) en une interface unifiée.
Le backend, développé avec Django REST Framework, expose une API REST robuste
qui orchestre les analyses de sécurité et gère l'authentification via GitHub OAuth~2.0.
Le frontend React/TypeScript offre une expérience utilisateur intuitive permettant
de déclencher des analyses, consulter les résultats et interagir avec l'assistant
IA. L'intégration d’un modèle de langage via Ollama, couplé à une approche RAG (Retrieval-Augmented Generation), pour la mise en place d’un système de recommandation des scanners de sécurité et des solutions de remédiation, représente une contribution originale qui rend la plateforme accessible même aux développeurs non spécialistes en sécurité. Enfin, l'intégration avec DefectDojo assure un suivi structuré et traçable des vulnérabilités identifiées.

Sur le plan des perspectives, plusieurs axes d'amélioration peuvent être envisagés
pour enrichir VulnOps à l'avenir :

\begin{enumerate}
    \item \textbf{Intégration CI/CD native :} L'ajout d'un plugin GitHub Actions
    permettrait de déclencher automatiquement les analyses VulnOps à chaque
    \textit{pull request}, renforçant ainsi l'intégration continue de la sécurité.
    
    \item \textbf{Déploiement cloud et SaaS :} VulnOps pourrait être transformé
    en service SaaS accessible depuis un navigateur web, avec une facturation
    à l'usage, rendant la plateforme accessible à des organisations de toutes
    tailles sans infrastructure locale.
    
    \item \textbf{Tableaux de bord analytiques :} L'ajout de visualisations
    temporelles permettrait de suivre l'évolution de la posture de sécurité
    d'un projet sur la durée et d'identifier les tendances dans les types
    de vulnérabilités détectées.
\end{enumerate}

%============================================================
% BIBLIOGRAPHIE
%============================================================
\begin{thebibliography}{99}

\bibitem{devsecopsgartner}
    Gartner Research,
    \textit{"DevSecOps: How to Seamlessly Integrate Security Into DevOps"},
    Gartner Inc., 2019.

\bibitem{owasptop10}
    OWASP Foundation,
    \textit{"OWASP Top Ten 2021 -- The Ten Most Critical Web Application Security Risks"},
    Open Web Application Security Project, 2021.
    Disponible sur : \url{https://owasp.org/www-project-top-ten/}

\bibitem{defectdojo}
    Greg Anderson, et al.,
    \textit{"DefectDojo -- The Open-Source Application Vulnerability Management Tool"},
    GitHub Repository, 2025.
    Disponible sur : \url{https://github.com/DefectDojo/django-DefectDojo}

\bibitem{oauth2rfc}
    D. Hardt (Ed.),
    \textit{"The OAuth 2.0 Authorization Framework"},
    RFC 6749, Internet Engineering Task Force (IETF), 2012.
    Disponible sur : \url{https://datatracker.ietf.org/doc/html/rfc6749}

\bibitem{umlspec}
    Object Management Group (OMG),
    \textit{"Unified Modeling Language (UML) Specification, Version 2.5.1"},
    OMG Document formal/17-12-05, 2017.

\bibitem{cleanarchitecture}
    Robert C. Martin,
    \textit{"Clean Architecture: A Craftsman's Guide to Software Structure and Design"},
    Prentice Hall, 2017.

\bibitem{gof}
    Erich Gamma, Richard Helm, Ralph Johnson, John Vlissides,
    \textit{"Design Patterns: Elements of Reusable Object-Oriented Software"},
    Addison-Wesley, 1994.

\bibitem{scrumguide}
    Ken Schwaber, Jeff Sutherland,
    \textit{"The Definitive Guide to Scrum: The Rules of the Game"},
    Scrum.org, 2020.

\bibitem{djangodocs}
    Django Software Foundation,
    \textit{"Django Documentation, Version 4.2"},
    2023.
    Disponible sur : \url{https://docs.djangoproject.com/}

\bibitem{reactdocs}
    Meta Open Source,
    \textit{"React -- A JavaScript library for building user interfaces"},
    React Documentation, 2024.
    Disponible sur : \url{https://react.dev/}

\bibitem{banditdoc}
    PyCQA,
    \textit{"Bandit -- A tool designed to find common security issues in Python code"},
    GitHub Repository, 2024.
    Disponible sur : \url{https://github.com/PyCQA/bandit}

\bibitem{semgrepdoc}
    Semgrep Inc.,
    \textit{"Semgrep -- Static analysis at ludicrous speed"},
    Semgrep Documentation, 2024.
    Disponible sur : \url{https://semgrep.dev/docs/}

\bibitem{ollama}
    Ollama,
    \textit{"Ollama -- Get up and running with large language models locally"},
    2024.
    Disponible sur : \url{https://ollama.com/}

\bibitem{gitpython}
    Sebastian Thiel,
    \textit{"GitPython -- Python library used to interact with Git repositories"},
    GitHub Repository, 2024.
    Disponible sur : \url{https://github.com/gitpython-developers/GitPython}

\bibitem{corssec}
    W3C,
    \textit{"Cross-Origin Resource Sharing (CORS)"},
    W3C Recommendation, 2014.
    Disponible sur : \url{https://www.w3.org/TR/cors/}

\end{thebibliography}

%============================================================
% NETOGRAPHIE
%============================================================
\chapter*{Netographie}
\addcontentsline{toc}{chapter}{Netographie}
\markboth{Netographie}{Netographie}

\begin{enumerate}
    \item \url{https://owasp.org/} -- Site officiel de l'OWASP Foundation, consulté en mars 2026.
    \item \url{https://docs.djangoproject.com/} -- Documentation officielle Django, consultée en février 2026.
    \item \url{https://react.dev/} -- Documentation officielle React, consultée en février 2026.
    \item \url{https://www.typescriptlang.org/docs/} -- Documentation TypeScript, consultée en mars 2026.
    \item \url{https://vitejs.dev/} -- Documentation Vite, consultée en mars 2026.
    \item \url{https://github.com/DefectDojo/django-DefectDojo} -- Dépôt GitHub de DefectDojo, consulté en avril 2026.
    \item \url{https://semgrep.dev/docs/} -- Documentation Semgrep, consultée en mars 2026.
    \item \url{https://github.com/PyCQA/bandit} -- Dépôt GitHub de Bandit, consulté en mars 2026.
    \item \url{https://ollama.com/library} -- Bibliothèque des modèles Ollama, consultée en mars 2026.
    \item \url{https://docs.github.com/en/apps/oauth-apps} -- Documentation GitHub OAuth Apps, consultée en février 2026.
    \item \url{https://ollama.com/} -- Site officiel d'Ollama (inférence LLM locale), consulté en avril 2026.
    \item \url{https://pylint.pycqa.org/} -- Documentation PyLint, consultée en mars 2026.
    \item \url{https://www.sonarcloud.io/} -- Plateforme SonarCloud, consultée en mars 2026.
    \item \url{https://datatracker.ietf.org/doc/html/rfc6749} -- RFC 6749 OAuth 2.0, consultée en février 2026.
    \item \url{https://www.docker.com/} -- Documentation Docker, consultée en avril 2026.
\end{enumerate}

%============================================================
% ANNEXE
%============================================================
\appendix
\chapter{État de l'art : Techniques d'analyse statique de sécurité (SAST)}
\label{annexe:sast}

\section{Présentation générale du SAST}

L'analyse statique de sécurité des applications (SAST, \textit{Static Application
Security Testing}) est une méthode d'examen du code source, du bytecode ou des
binaires d'une application afin d'identifier des vulnérabilités de sécurité
potentielles, et ce, sans exécu\-ter le programme \cite{owasptop10}. Contrairement
à l'analyse dynamique (DAST), le SAST intervient en amont du déploiement,
permettant une détection précoce des failles, réduisant ainsi le coût de correction.

\section{Catégories de vulnérabilités détectées}

Les outils SAST sont capables de détecter plusieurs catégories de vulnérabilités,
notamment :

\begin{itemize}
    \item \textbf{Injection (SQL, commandes OS, LDAP) :} Utilisation non sécurisée
    de données externes dans des requêtes ou commandes.
    \item \textbf{Cryptographie faible :} Utilisation d'algorithmes de hachage
    obsolètes (MD5, SHA-1) ou de clés de taille insuffisante.
    \item \textbf{Gestion des erreurs inadéquate :} Exposition d'informations
    sensibles dans les messages d'erreur.
    \item \textbf{Secrets codés en dur :} Mots de passe, clés API ou tokens
    présents directement dans le code source.
    \item \textbf{Cross-Site Scripting (XSS) :} Absence de validation et
    d'échappement des données utilisateur dans les sorties HTML.
\end{itemize}

\section{Présentation des outils intégrés dans VulnOps}

\subsection{Bandit}

Bandit \cite{banditdoc} est un outil SAST open-source développé par PyCQA
spécialisé dans l'analyse du code Python. Il examine l'arbre syntaxique abstrait
(AST) du code et détecte des patterns de code dangereux en utilisant un système
de plugins. Bandit génère des rapports dans de multiples formats (JSON, HTML, CSV)
et classe les vulnérabilités selon leur sévérité (High, Medium, Low) et leur
niveau de confiance.

\subsection{Semgrep}

Semgrep \cite{semgrepdoc} est un outil d'analyse statique multi-langages
basé sur la correspondance de patterns syntaxiques. Il supporte plus de vingt
langages de programmation et dispose d'un registre de règles communautaires
couvrant les vulnérabilités OWASP Top 10. Dans VulnOps, Semgrep est configuré
avec la règle \texttt{p/owasp-top-ten} pour une couverture maximale des
vulnérabilités web courantes.

\subsection{ESLint (mode sécurité)}

ESLint est principalement connu comme un linter JavaScript/TypeScript, mais
dispose d'un ensemble de règles dédiées à la sécurité applicative via le
plugin \texttt{eslint-plugin-security}. Ces règles permettent de détecter
des usages dangereux comme l'évaluation de code dynamique (\texttt{eval()}),
les injections d'expressions régulières (ReDoS), et les lectures de fichiers
non sécurisées.

\chapter{Guide de déploiement de VulnOps}
\label{annexe:deploiement}

\section{Prérequis}

Pour déployer VulnOps, les pré-requis suivants sont nécessaires :
Python 3.11+, Node.js 20+, Git, Docker et Docker Compose (pour DefectDojo).

\section{Installation du backend}

\begin{lstlisting}[language=bash, caption=Installation du backend Django]
# Créer l'environnement virtuel Python
python -m venv venv
source venv/bin/activate  # Windows: venv\Scripts\activate

# Installer les dépendances
pip install -r backend/requirements.txt

# Configurer les variables d'environnement
cp .env.example .env
# Editer .env avec vos clés API

# Appliquer les migrations
cd backend
python manage.py migrate

# Lancer le serveur de développement
python manage.py runserver
\end{lstlisting}

\section{Installation du frontend}

\begin{lstlisting}[language=bash, caption=Installation du frontend React]
cd frontend
npm install
npm run dev
\end{lstlisting}

\section{Configuration des variables d'environnement}

\begin{lstlisting}[caption=Fichier .env principal]
# GitHub OAuth
GITHUB_CLIENT_ID=your_github_client_id
GITHUB_CLIENT_SECRET=your_github_client_secret

# Ollama (Local LLM)
OLLAMA_BASE_URL=http://localhost:11434
OLLAMA_MODEL=mistral

# DefectDojo
DEFECTDOJO_URL=http://localhost:8080
DEFECTDOJO_API_KEY=your_defectdojo_api_key

# Django
SECRET_KEY=your_django_secret_key
DEBUG=True
ALLOWED_HOSTS=localhost,127.0.0.1
\end{lstlisting}

%============================================================
% FIN DU DOCUMENT
%============================================================
\begin{center}
    \vspace{3cm}
    \textcolor{primaryblue}{\rule{4cm}{1.5pt} \quad \textbf{FIN DU RAPPORT} \quad \rule{4cm}{1.5pt}}
\end{center}

\end{document}
